\chapter{Background and related work}\label{chap:background}

\section{The Loss of Perimeter Security}\label{sec:loss-of-perimeter-security}
The  security  of modern networks has been built around a well-defined perimeter, a topological barrier that separates trusted local networks (LAN) from wider, potentially hostile networks (WAN). However, the transition toward IPv6 alters this defensive boundary by removing legacy mechanisms and creating critical asymmetries in firewall policies. For decades, NAT protected IoT devices by defaulting to deny inbound traffic. However, with IPv6's direct global addressing, this implicit security disappears. As the research below demonstrates, organizations have failed to compensate with explicit firewall policies, creating a ``perfect storm'' of exposure~\cite{rye2025firewalls, czyz2016backdoor}.

The elimination of NAT creates the need for IPv6 networks to implement filtering policies and firewalls that are equivalent to, or stronger than, those used in IPv4. Nonetheless, empirical research shows a substantial gap in the enforcement of these organizational policies. A large-scale characterization of dual-stack networks analyzing the behavior of 520,000 servers and 25,000 routers demonstrated that corporate security policies in IPv6 are consistently more permissive~\cite{czyz2016backdoor}. In particular, several high-value target applications were identified with more permissive security policies in IPv6. Management protocols commonly targeted by attackers and botnets, such as SSH, Telnet, and SNMP, were found to be more than twice as exposed on routers in IPv6 as they are in IPv4~\cite{czyz2016backdoor}. This systemic weakness suggests that organizations are more inclined to prioritize securing legacy infrastructure while leaving IPv6 deployments insufficiently protected due to permissive configurations or technical and operational oversight~\cite{czyz2016backdoor}.

While corporate networks already show significant security gaps, the situation in residential networks is more concerning. Recent studies on residential IPv6 deployments highlight significant security risks related with the residential transitions. Differently from IPv4 NAT, which by default imposes a default-deny access control model, IPv6 residential router manufacturers (CPE) have the discretion to decide whether to implement restrictive or permissive filtering policies. After actively scanning more than 2.5 million residential IPv6 prefixes, the researchers found that many routers adopted a default-allow policy~\cite{rye2025firewalls}, deliberately routing unsolicited incoming traffic from the WAN to the LAN. The study further showed that simple IPv6 scans allowed researchers to directly reach end devices such as printers, and smart light bulbs, surpassing the discovery rates that previously required exhaustive scanning of the entire IPv4 address space. These findings indicate that the residential transition weakens traditional firewall protections, exposing additional hosts to botnets and increasing overall Internet security risks~\cite{rye2025firewalls}.

\section{The IoT Security Landscape}\label{sec:iot-landscape}
The loss of the network perimeter described above conflicts with the current state of the embedded device ecosystem (IoT). The design of these devices often assumes that their operating environment will remain isolated and private, resulting in firmware architectures and communication protocols that are vulnerable when directly exposed to the Internet through IPv6.

Security practices in home IoT devices have been extensively evaluated~\cite{alrawi2019sok}. After evaluating 45 commercial devices across four components: the device itself, the mobile application, the cloud endpoint, and the communication channel, this evaluation revealed that manufacturers often assume that the LAN environment is inherently trustworthy and protected by a boundary router. As a result, many devices rely primarily on the domestic router as the main security mechanism. This assumption appears as a consequence of prioritizing product release time over security considerations, which leads to the omission of internal security controls. Many devices lack local encryption mechanisms, making them susceptible to man-in-the-middle attacks, while mobile applications frequently request excessive privileges that go beyond their intended functionality~\cite{alrawi2019sok}. These results are not limited to unbranded IoT devices. Platforms such as Samsung SmartThings also exhibit inherent design flaws in their permission model that result in over-privileged third-party applications, allowing potential remote exploitation when the devices are externally accessible~\cite{fernandes2016smarthome}. Considering these findings together with the disappearance of perimeter filtering policies, the implication is that devices and platforms designed under the assumption of network isolation become directly exposed to remote attackers.

In the application layer, the technical debt of the embedded ecosystem increases the viability of exploitation. Hardware limitations and accelerated development cycles lead to devices running obsolete or misconfigured network services. An examination of 32,000 firmware images discovered 38 previously unknown vulnerabilities in more than 693 firmwares, confirming that latent vulnerabilities persist due to outdated software and the fact that security is systematically sacrificed for cost and time-to-market~\cite{costin2014firmware}. The results are consistent with broader evaluations of domestic IoT devices~\cite{alrawi2019sok}, which revealed a persistent reliance on plaintext protocols such as HTTP instead of implementing TLS encryption, as well as the use of legacy administrative services such as Telnet. In traditional IPv4 networks behind NAT, interacting with these vulnerabilities would require prior access to the LAN. However, in the context of globally routable IPv6 networks, these insecure services can be directly exposed, facilitating credential extraction and the propagation of automated malware. The Mirai botnet demonstrated the consequences of this type of exposure, using a dictionary of only 62 pairs of default credentials for brute-force attacks via Telnet, the malware was able to infect more than 600,000 IoT devices and execute massive DDoS attacks against high-profile targets~\cite{antonakakis2017mirai}.

Even when manufacturers protect public web interfaces with passwords, devices may still expose residual risk through hidden administrative channels, as documented in routing analysis assessments~\cite{liu2025eagleye}. Hidden network interfaces are defined as functional access routes to the device that are neither documented nor accessible through the standard navigation of the main portal, often introduced during development but later abandoned in production~\cite{liu2025eagleye}. Because they are not exposed through the usual interaction flow, these interfaces evade standard authentication controls. A systematic analysis of 13 commercial devices discovered 79 hidden interfaces, 25 times more than previous tools commonly used. The activation of these hidden channels triggered 29 zero-day vulnerabilities, including command injection, XSS, information leakage, and operating backdoors that provided full access to the device, resulting in seven assigned CVEs~\cite{liu2025eagleye}. The presence of this type of authentication-bypass vulnerability in embedded firmware has also been validated using binary analysis frameworks such as Firmalice, which detected backdoors that allowed administrative access through specially crafted packets without requiring legitimate credentials~\cite{shoshitaishvili2015firmalice}. When these findings are considered together, the erosion of the local network perimeter not only exposes known vulnerabilities in IoT devices but also opens the path for potential system takeover and credential bypass.

\section{IPv6 Visibility \& Identifiability}\label{sec:ipv6-visibility}
Even considering the scenario described above, a defensive argument that justified several of the shortcomings of the IoT ecosystem was that, regardless of its vulnerabilities, the 128-bit IPv6 address space provided a form of practical invisibility, in a role similar to NAT in IPv4, different in mechanism but often assumed to produce a comparable protective effect~\cite{williams20246sense}. Nonetheless, core properties of IPv6, particularly its reliance on Stateless Address Autoconfiguration (SLAAC) without mandatory cryptographic authentication~\cite{rfc8200,rfc3756} and the deterministic embedding of hardware MAC addresses via EUI-64~\cite{rye2023ipvseeyou}, make devices highly identifiable and enable the construction of precise visual and topological profiles for attackers.

In order to guarantee plug-and-play functionality, IPv6~\cite{rfc8200} relies on Stateless Address Autoconfiguration (SLAAC)~\cite{rfc4862}, a mechanism that uses mandatory messages in the Neighbor Discovery Protocol (NDP)~\cite{rfc4861}, specifically Router Solicitations (RS) and Router Advertisements (RA). The structural weakness lies in the fact that the IPv6 standard does not require cryptographic authentication for these local network exchanges. The threats resulting from this absence were formally cataloged by the IETF~\cite{rfc3756}, which documented how non-privileged actors within the same network segment can orchestrate man-in-the-middle attacks through falsified NDP messages, sending fake RA packets to reroute traffic or forcing denial of service during the autoconfiguration process. Even though Secure Neighbor Discovery (SEND)~\cite{rfc3971} was proposed as a mitigation mechanism, its implementation complexity and its architectural requirement for public keys have limited its practical deployment. The issue is further complicated because indiscriminately blocking ICMPv6~\cite{rfc4443} would break network connectivity, since NDP depends on these messages to operate~\cite{rfc4890}. As a result, the network initialization mechanisms used by IPv6, and therefore by many IoT devices, represent a permanently exposed attack surface.

The SLAAC process requires that the device generates its own interface identifier (IID) to use the last 64 bits of the IPv6 address. Traditionally, the industry has implemented the EUI-64 standard, which directly embeds the physical address (MAC) of the hardware into the public IID. The privacy and security implications of this design have been demonstrated through empirical measurements of EUI-64 address leakage~\cite{rye2023ipvseeyou}. These measurements revealed that this mechanism produces a systematic leakage of MAC addresses, allowing an unprivileged remote adversary to geolocate residential routers with street-level precision, identifying more than 12 million routers in 146 countries and territories with a median geolocation error of 39 meters~\cite{rye2023ipvseeyou}. By correlating the Organizationally Unique Identifier (OUI) with IPv6 geolocation databases and Wi-Fi positioning data, the investigation showed that an identifier at the network layer had become a direct vector for physical tracking and the identification of the device manufacturer. Once these identification capabilities are combined with firmware vulnerabilities (Section~\ref{sec:iot-landscape}), a remote attacker could not only discover the IoT device but also determine the most likely manufacturer, firmware, and physical location, enabling highly targeted attacks.

The visibility of the connected ecosystem is further increased by the traffic generated at the application layer that modern protocols require to function properly. IoT devices use multicast protocols such as mDNS and SSDP for service discovery within the local network~\cite{yu2020broadcast}. This passive traffic has been shown to form a highly descriptive network fingerprint~\cite{yu2020broadcast}. Through the analysis of structural and textual information embedded in broadcast and multicast packets, the proposed framework was able to identify with high accuracy the brand, model, and category of IoT and mobile devices without actively interacting with them~\cite{yu2020broadcast}. In an IPv6 environment where devices are directly reachable, this passive profiling capability complements manufacturer identification through EUI-64~\cite{rye2023ipvseeyou}, giving an attacker a detailed map of the device ecosystem before launching any exploit.

Summarizing, these findings destroy the illusion of opacity. In an IPv6 environment, an IoT device not only receives a publicly reachable address by default, but its autoconfiguration mechanisms also expose the identity of its manufacturer through EUI-64 addresses~\cite{rye2023ipvseeyou}. At the same time, NDP exchanges do not require mandatory cryptographic authentication~\cite{rfc3756}, and broadcast protocols can passively reveal the device’s brand, model, and category to any listener on the network segment~\cite{yu2020broadcast}. As a result, the assumed ``invisibility'' of IoT devices disappears by design in the IPv6 environment.

\section{Scanning \& Target Generation in IPv6}\label{sec:scanning-ipv6}
The scanning and discovery phase is the first step in the chain of a cyberattack. Under IPv4, with an address space of only $2^{32}$ addresses, tools such as ZMap showed that it was possible to scan the entire Internet in less than one hour~\cite{durumeric2013zmap}. IPv6, with its much larger address space, was designed under the assumption that its virtually unlimited size would make exhaustive enumeration computationally impractical, turning address space size into a form of implicit security. However, recent studies have shown that this mathematical barrier has been systematically overcome through smart and adaptive scanning techniques.

The capability to scan the entire IPv4 space expanded large-scale Internet measurement. Building on tools such as ZMap~\cite{durumeric2013zmap}, platforms like Censys allowed the indexing of vulnerabilities, certificates, and embedded devices across the Internet~\cite{durumeric2015censys}. This visibility over the IPv4 address space made it possible to identify services and devices exposed at a global scale. Botnets such as Mirai exploited the same principle, locating IoT devices such as IP cameras by performing simple stochastic scans and compromising them through default Telnet credentials, as discussed in Section~\ref{sec:iot-landscape}.

The IPv6 address space was designed under the idea that it would be too large to make exhaustive mapping computationally feasible. Nevertheless, recent research has shown that this space can still be mapped through adaptive scanning systems. One of the most notable contributions is \textit{6Sense}~\cite{williams20246sense}. One of the main challenges in IPv6 scanning is aliasing, where intermediate routers that are poorly configured respond to scans directed at tens of millions of inactive addresses within a prefix, diluting the resources of the scanner with false positives~\cite{williams20246sense}. 6Sense addresses this problem through a reinforcement learning approach combined with online and offline de-aliasing modules that can distinguish between illusory prefixes~\cite{williams20246sense}. By dynamically learning active addresses and focusing on predictable assignment patterns—such as lower-order bytes used by IoT interfaces or grouping patterns in nybbles—the system was able to discover tens of millions of active IPv6 hosts, up to $3.6\times$ more active hosts than previous methods, and successfully correlate them with critical vulnerabilities (CVEs)~\cite{williams20246sense}. Furthermore, side channels derived from ICMP rate limiting in routers allow the structure of IPv6 networks to be inferred without the need for direct scans~\cite{pan2023router}. These results establish that the virtually unlimited IPv6 address space is not a viable mechanism for obscurity.


\section{Synthesis \& Research Gap}\label{sec:synthesis-gap}
The analysis of the literature and related work presented in the previous sections reveals potential issues associated with the transition to IPv6. The previous sections examined four dimensions of this problem individually: the disappearance of the network perimeter within which the IoT environment was originally developed (Section~\ref{sec:loss-of-perimeter-security}), the inherent vulnerabilities of IoT devices (Section~\ref{sec:iot-landscape}), the identifiability of devices in IPv6 networks (Section~\ref{sec:ipv6-visibility}), and the feasibility of large-scale scanning in the IPv6 address space (Section~\ref{sec:scanning-ipv6}). When considered in isolation, each of these phenomena represents a considerable risk. However, when examined together, a qualitatively different picture emerges. Despite this body of work, a gap in current research methodologies remains.

In a complete IPv6 environment, an IoT device connected to a residential network automatically receives a globally routable address, and in many cases the router does not filter incoming traffic~\cite{rye2025firewalls, rfc4890}. At the moment the device is autoconfigured through SLAAC, which relies on mechanisms that lack cryptographic authentication~\cite{rfc3756}, the identity of the manufacturer may already be exposed through EUI-64~\cite{rye2023ipvseeyou}, while its discovery protocols may have already passively revealed its brand, model, and category~\cite{yu2020broadcast}. At the same time, specialized scanning tools such as 6Sense can locate active devices within the 128-bit IPv6 address space~\cite{williams20246sense}. Once discovered and identified, the device presents the same vulnerabilities extensively documented in the literature: firmware that runs unsafe protocols or lacks encryption~\cite{alrawi2019sok}, default credentials that can be exploited to obtain administrator privileges as demonstrated by the Mirai botnet~\cite{antonakakis2017mirai}, hidden interfaces that bypass authentication~\cite{liu2025eagleye}, and platforms with deficient permission hierarchies~\cite{fernandes2016smarthome}.

Each link of this chain has been individually validated, but the full chain remains unexplored. Network-focused studies concentrate on scanning feasibility and filtering policies without evaluating what happens once a device is reached. Device-focused research analyzes firmware vulnerabilities assuming an isolated IPv4 local network. No existing study examines, under controlled conditions, which specific router configurations and device settings enable or prevent the external reachability of an IoT device in a residential IPv6 environment.

This thesis addresses this gap through a controlled laboratory experiment using commercial hardware (ESP32), simulating being an IoT device, connected to real residential IPv6 infrastructure. By systematically varying router configurations and device parameters, this work identifies the specific conditions under which an IoT device with a globally routable IPv6 address becomes externally accessible, and demonstrates the security implications of that exposure.
